\section{Структура HeightHandler}

Ключевым компонентом решения является структура \texttt{HeightHandler}, которая эффективно хранит информацию о текущих высотах на паллете. Данная структура критически важна для производительности алгоритма, поскольку операции с ней выполняются при размещении каждой коробки.

\subsection{Основные операции}

HeightHandler предоставляет три ключевые операции:

\begin{itemize}
    \item \texttt{get\_h(x, y, X, Y)} --- возвращает максимальную высоту в прямоугольной области $[x, X] \times [y, Y]$
    \item \texttt{get\_area(x, y, X, Y)} --- возвращает площадь (в мм²), которая находится на максимальной высоте в заданном прямоугольнике (используется для расчёта опоры коробки)
    \item \texttt{add\_rect(x, y, X, Y, h)} --- добавляет новый прямоугольник с заданной высотой, обновляя карту высот
    \item \texttt{get\_dots(header, box)} --- возвращает ``интересные'' точки для размещения коробки заданного размера
\end{itemize}

Функция \texttt{get\_dots} генерирует кандидатные позиции на основе углов существующих прямоугольников, что значительно сокращает пространство поиска с $O(W \times L)$ до $O(n)$, где $n$ --- количество уже уложенных коробок.

% TODO: Добавить схему работы HeightHandler (figures/height_handler_scheme.png)
% Показать: 2D-вид паллеты сверху с прямоугольниками разных высот (цветовая карта)
% и точки get_dots вокруг углов прямоугольников

\subsection{Реализации HeightHandler}

В рамках работы были исследованы пять различных реализаций структуры HeightHandler с разными алгоритмическими подходами:

\subsubsection{Baseline (наивная реализация)}

Простейшая реализация, которая хранит все добавленные прямоугольники в векторе и при каждом запросе честно перебирает их.

\textbf{Сложность операций:}
\begin{itemize}
    \item \texttt{get\_h()}: $O(n)$ --- перебор всех прямоугольников
    \item \texttt{get\_area()}: $O(n)$ --- перебор всех прямоугольников
    \item \texttt{add\_rect()}: $O(1)$ --- просто добавление в конец вектора
\end{itemize}

\textbf{Особенность:} Корректность гарантируется свойством задачи паллетизации --- прямоугольники добавляются ``сверху'', поэтому новый прямоугольник всегда имеет высоту строго больше, чем все существующие в той же области. Два прямоугольника на одной высоте никогда не перекрываются.

\subsubsection{Rects (оптимизированный список прямоугольников)}

Улучшенная версия Baseline с оптимизацией: прямоугольники сортируются по убыванию высоты. Это позволяет при поиске максимума остановиться на первом найденном пересечении.

\textbf{Сложность операций:}
\begin{itemize}
    \item \texttt{get\_h()}: $O(n)$ в худшем случае, но на практике значительно быстрее из-за раннего выхода
    \item \texttt{get\_area()}: $O(n)$ в худшем случае, раннее завершение при снижении высоты
    \item \texttt{add\_rect()}: $O(n)$ --- вставка с сохранением сортировки
\end{itemize}

\subsubsection{Grid (плотная 2D сетка)}

Разбивает паллету на ячейки размером $10 \times 10$ мм. В каждой ячейке хранится максимальная высота. Для паллеты $1200 \times 800$ мм это $120 \times 80 = 9600$ ячеек.

\textbf{Сложность операций:}
\begin{itemize}
    \item \texttt{get\_h()}: $O(k)$, где $k$ --- количество ячеек в запросе
    \item \texttt{get\_area()}: $O(k)$
    \item \texttt{add\_rect()}: $O(k)$ --- обновление всех затронутых ячеек
\end{itemize}

\textbf{Недостаток:} Большой объём памяти и медленное обновление при добавлении крупных прямоугольников.

\subsubsection{SegTree (2D дерево отрезков)}

Двумерное дерево отрезков с отложенным обновлением (Lazy Propagation). X-дерево содержит Y-деревья, каждый Y-узел хранит количество 2D-ячеек с максимальной высотой.

\textbf{Сложность операций:}
\begin{itemize}
    \item \texttt{get\_h()}: $O(\log W \cdot \log H)$
    \item \texttt{get\_area()}: $O(\log W \cdot \log H)$
    \item \texttt{add\_rect()}: $O(\log W \cdot H)$ --- необходим merge Y-деревьев
\end{itemize}

\textbf{Особенность:} Асимптотически эффективен для запросов, но высокая константа и сложная операция обновления делают его медленнее на практике.

\subsubsection{Quadtree (квадродерево)}

Рекурсивное разбиение пространства на 4 части. Каждый узел хранит максимальную высоту в своей области.

\textbf{Сложность операций:}
\begin{itemize}
    \item \texttt{get\_h()}: $O(\log(\max(W,H)))$ при равномерном распределении
    \item \texttt{get\_area()}: $O(n)$ --- использует fallback на список прямоугольников
    \item \texttt{add\_rect()}: $O(\log(\max(W,H)))$ в среднем, но может деградировать
\end{itemize}

\textbf{Недостаток:} Высокие накладные расходы на управление памятью (unique\_ptr), частое разбиение узлов.

\subsection{Сравнительное тестирование}

Тестирование проводилось на 436 реальных задачах с суммарно 27.5 миллионами операций.

\begin{table}[ht]
\centering
\caption{Результаты тестирования реализаций HeightHandler (время в мс)}
\label{tab:heighthandler}
\begin{tabular}{|l|r|r|r|r|r|}
\hline
\textbf{Алгоритм} & \textbf{Total} & \textbf{add\_rect} & \textbf{get\_h} & \textbf{get\_area} & \textbf{Speedup} \\
\hline
Baseline & 10010.8 & 2.2 & 4471.9 & 5706.6 & 1.00x \\
Rects & 775.1 & 9.1 & 551.4 & 671.8 & \textbf{12.92x} \\
Grid & 30491.4 & 33.0 & 6414.5 & 24322.6 & 0.33x \\
SegTree & 19749.0 & 1774.8 & 9795.1 & 8606.3 & 0.51x \\
Quadtree & 56304.0 & 325.1 & 30599.0 & 26120.6 & 0.18x \\
\hline
\end{tabular}
\end{table}

\subsection{Анализ результатов}

\textbf{Rects} показал наилучшую производительность с ускорением в \textbf{12.92x} относительно базовой реализации. Это объясняется:
\begin{itemize}
    \item Сортировка по убыванию высоты позволяет быстро находить максимум
    \item Низкие накладные расходы (простой вектор)
    \item Эффективное использование кэша процессора
\end{itemize}

\textbf{Grid, SegTree и Quadtree} оказались медленнее Baseline. Причины:
\begin{itemize}
    \item Высокие накладные расходы на управление структурами данных
    \item Типичные коробки покрывают значительную часть паллеты, что нивелирует преимущества иерархических структур
    \item Константные факторы асимптотически эффективных алгоритмов превышают выигрыш
\end{itemize}

В финальной версии алгоритма используется реализация \textbf{Rects}. Вот так, несмотря на простую реализацию, она выигрывает у более мощных алгоритмов.
